\chapter{Conclusão}

A presente prática laboratorial cumpriu plenamente seus objetivos ao caracterizar, de forma teórica e experimental, a dinâmica de filtros ativos passa-baixa de segunda ordem baseados na topologia de Múltipla Realimentação (MFB). A correlação entre os modelos analíticos desenvolvidos no domínio de Laplace e os dados empíricos extraídos em bancada confirmou a validade e a robustez da aproximação do amplificador operacional ideal para o projeto de circuitos lineares.

O principal mérito do experimento residiu na observação direta do impacto do fator de amortecimento ($\zeta$) na resposta do sistema. Ao permutar os resistores $R_1$ e $R_3$, evidenciou-se a drástica transição na dinâmica do circuito:
\begin{itemize}
	\item \textbf{No Primeiro Circuito ($\zeta > 1$):} Validou-se o regime superamortecido, caracterizado por uma atenuação monotônica e suave da magnitude, com uma frequência de corte real mensurada em $4,3\text{ Hz}$ e ausência total de sobressinal na resposta temporal.
	\item \textbf{No Segundo Circuito ($0 < \zeta < 1$):} Comprovou-se o regime subamortecido, evidenciado pelo surgimento de um pico de ressonância no diagrama de Bode exatamente em $17,5\text{ Hz}$, seguido de uma atenuação abrupta na banda de rejeição, comportamento típico de sistemas com polos complexos conjugados.
\end{itemize}

A metodologia de aferir os valores nominais a frio demonstrou-se essencial. A inserção das resistências e capacitâncias reais nas funções de transferência recalculadas justificou os pequenos desvios de frequência encontrados na bancada, provando que os erros não decorreram de falhas de montagem ou limites do amplificador operacional, mas sim da tolerância comercial dos componentes passivos. 

Em suma, a prática consolidou competências fundamentais em instrumentação eletrônica e análise de sistemas de segunda ordem. A compreensão profunda de como dimensionar os polos de um circuito para evitar oscilações indesejadas (ou para forçar ressonâncias seletivas) constitui um pilar essencial para o futuro projeto de sistemas de condicionamento de sinais, telecomunicações e controle.